\begin{table}[htbp]
\centering
\caption{Perbandingan alur perawatan pada kondisi saat ini terhadap rancangan yang diusulkan.}
\label{tab:before-after}
\begin{tabular}{@{}>{\raggedright\arraybackslash\hyphenpenalty=10000}p{2.6cm}>{\raggedright\arraybackslash\hyphenpenalty=10000}p{4.4cm}>{\raggedright\arraybackslash\hyphenpenalty=10000}p{4.4cm}c@{}}
\toprule
\textbf{Aspek} & \textbf{Kondisi saat ini} & \textbf{Rancangan yang diusulkan} & \textbf{Baru} \\
\midrule
Sifat tindakan & Reaktif; tindakan diambil sesudah penyimpangan teramati & Antisipatif; panduan disiapkan bagi keadaan yang diperkirakan akan terjadi & \\[3pt]
Penafsiran catatan & Manual oleh dokter pada saat konsultasi & Kondisi klinis terstruktur berupa tingkat risiko, arah perubahan, dan kegentingan & \\[3pt]
Sumber kondisi penyusun kueri & Tidak ada penelusuran otomatis & Kondisi yang \textbf{diprediksi}, bukan kondisi yang sedang berlaku & $\bigstar$ \\[3pt]
Panduan klinis & Dicari manual pada dokumen, tanpa catatan asal & Ditelusuri otomatis, disertai penunjuk dokumen dan nomor halaman & \\[3pt]
Peran dokter & Penentu tunggal tanpa bahan pendukung terstruktur & Penentu akhir atas rekomendasi yang dapat diperiksa asalnya & \\
\bottomrule
\end{tabular}

\vspace{4pt}
\noindent\tabnotestyle Tanda $\bigstar$ menandai baris yang menjadi letak kebaruan penelitian ini. Keempat baris lainnya merupakan konsekuensi perancangan yang lazim dijumpai pada sistem pendukung keputusan klinis dan tidak diklaim sebagai kebaruan.\normalsize
\end{table}
